\documentclass[12pt,a4paper]{article}

\usepackage[utf8]{inputenc}
\usepackage[T1]{fontenc}
\usepackage[spanish]{babel}
\usepackage[margin=2.5cm]{geometry}
\usepackage{amsmath}
\usepackage{graphicx}
\usepackage{hyperref}
\usepackage{parskip}
\usepackage{listings}
\usepackage{xcolor}
\usepackage{booktabs}
\usepackage{tikz}
\usepackage{enumitem}
\usetikzlibrary{shapes.geometric, arrows.meta, positioning}

\hypersetup{colorlinks=true, linkcolor=blue!60!black, urlcolor=blue!60!black}

\lstset{
    basicstyle=\ttfamily\small,
    keywordstyle=\color{blue!70!black}\bfseries,
    commentstyle=\color{gray},
    stringstyle=\color{green!50!black},
    backgroundcolor=\color{gray!8},
    frame=single,
    rulecolor=\color{gray!40},
    breaklines=true,
    numbers=left,
    numberstyle=\tiny\color{gray},
    xleftmargin=1.5em,
}

% ── portada ──────────────────────────────────────────────────────────────────
\title{
    \LARGE \textbf{Proyecto 3}\\[0.4em]
    \large Universidad del Valle de Guatemala -- Bases de Datos 2\\
}
\author{
    Hugo Barillas \and
    Pablo Barillas \and
    Osman De León \and
    José Mérida \and
    Roberto Nájera \and
    Milton Polanco
}
\date{\today}

\begin{document}

\maketitle
\tableofcontents
\newpage

% ─────────────────────────────────────────────────────────────────────────────
\section{Descripción del Problema}

Se nos entregó un rompecabezas físico con piezas interconectadas mediante
\emph{tabs} (picos) y \emph{huecos}. El objetivo es construir un sistema que
almacene la estructura del rompecabezas en una base de datos y genere, a partir
de una pieza inicial, instrucciones ordenadas de ensamblaje. El sistema debe
manejar el caso de piezas faltantes y seguir armando la mayor parte posible del
rompecabezas.

Lo importante a notar es que el problema es fundamentalmente \emph{relacional}:
no importa la forma ni el color de cada pieza, sino únicamente qué conector de
qué pieza encaja con qué conector de otra pieza. Esta observación es la que
guía directamente la elección de tecnología.

% ─────────────────────────────────────────────────────────────────────────────
\section{Justificación de la Elección de Base de Datos}

\subsection{Tecnología seleccionada: Neo4j}

Se eligió \textbf{Neo4j} (base de datos orientada a grafos) porque refleja
directamente la estructura del problema. En un rompecabezas, el dato central no
es un atributo de una entidad aislada, sino la \emph{conexión} entre dos
conectores de piezas distintas. En Neo4j esa conexión es un arco de primer
nivel que existe por sí mismo; en un modelo relacional sería una fila en una
tabla de unión que solo cobra sentido al hacer un \textit{JOIN}. Eliminar ese
nivel de indirección simplifica tanto el modelo como las consultas.

Adicionalmente, Neo4j ofrece ventajas concretas para este dominio:

\begin{itemize}
    \item \textbf{Traversals nativos.} La pregunta ``¿qué pieza se conecta a
          esta pieza?'' es un salto de dos arcos en Cypher y no requiere lógica
          adicional. En SQL la misma consulta involucraría al menos cuatro tablas
          y varios \texttt{JOIN}.

    \item \textbf{Sin migraciones de esquema.} Añadir un nuevo rompecabezas es
          insertar un subgrafo aislado; no se modifica ninguna tabla existente ni
          se corre ningún script de migración.

    \item \textbf{Restricciones de unicidad.} Neo4j soporta restricciones
          \texttt{UNIQUE} equivalentes a claves primarias relacionales, lo que
          garantiza integridad referencial sin costo de modelado extra.
\end{itemize}

\subsection{Comparación con alternativas}

\textbf{PostgreSQL} hubiera sido una opción técnicamente válida, pero introduce
fricción innecesaria para este dominio. El modelo requeriría al menos cuatro
tablas (\texttt{puzzles}, \texttt{pieces}, \texttt{connectors},
\texttt{connections}) y cada consulta de vecindad involucraría múltiples
\texttt{JOIN}. El BFS en sí requeriría CTEs recursivos, trasladando la lógica
de traversal al SQL en lugar de mantenerla limpiamente en la capa de aplicación.
Para un problema donde la relación \emph{es} el dato, un modelo relacional
añade indirección sin ningún beneficio.

\textbf{MongoDB} maneja bien documentos anidados, lo que lo hace atractivo a
primera vista: una pieza con sus conectores encaja naturalmente en un documento.
El problema aparece con las conexiones entre piezas, que son simétricas: el
conector A de la pieza 1 y el conector B de la pieza 2 se referencian
mutuamente. Representar esa bidireccionalidad en documentos requiere duplicar la
referencia en ambos lados y mantenerlos sincronizados manualmente, introduciendo
una fuente de inconsistencia que Neo4j evita por diseño.

\textbf{DynamoDB} hubiera funcionado para rompecabezas muy pequeños con un
patrón de acceso simple y bien definido de antemano. Sin embargo, DynamoDB
está optimizado para lecturas y escrituras de ítems individuales por clave; las
consultas de vecindad que el BFS necesita requieren o bien un diseño de
\textit{access patterns} muy cuidadoso con índices secundarios, o bien múltiples
roundtrips a la base de datos por cada paso del algoritmo. A medida que el
rompecabezas crece, ese costo se vuelve prohibitivo. Para puzzles de juguete de
pocas piezas sería suficiente, pero no escala bien con la complejidad del grafo.

% ─────────────────────────────────────────────────────────────────────────────
\section{Diseño del Modelo de Datos}

\subsection{Convención de etiquetado físico}

Antes de cargar el rompecabezas al sistema, cada pieza y cada conector deben
etiquetarse físicamente para que el ensamblador pueda identificarlos a mano
durante el armado.

\begin{itemize}
    \item \textbf{Piezas}: se escribe un número en la parte delantera de cada
          pieza para identificarla de forma única dentro del rompecabezas.
          Se recomienda usar el formato \texttt{<número> - <nombre>} (p.ej.\
          \texttt{1 - Cabeza}) para que la identificación sea inequívoca incluso
          cuando varias piezas tienen formas similares.
    \item \textbf{Conectores}: se escribe una letra o código corto con cinta
          adhesiva junto a cada conector. La etiqueta es
          \emph{local} a su pieza, de modo que dos piezas distintas pueden tener
          un conector \texttt{A} sin conflicto.
\end{itemize}

Con este esquema, un conector queda identificado unívocamente por el par
$(\text{pieza},\ \text{etiqueta})$, p.ej.\ \texttt{(1, A)}, y una conexión
física queda registrada como el par de conectores que encajan entre sí.

\subsection{Nodos}

\begin{description}[leftmargin=3cm, style=sameline]
    \item[\texttt{:Puzzle}]    \texttt{id}, \texttt{name}, \texttt{description}
    \item[\texttt{:Piece}]     \texttt{id} (global), \texttt{local\_id}, \texttt{label}, \texttt{puzzle\_id}
    \item[\texttt{:Connector}] \texttt{id} (global), \texttt{local\_id}, \texttt{label}, \texttt{type} (\texttt{tab}/\texttt{hole}), \texttt{piece\_local\_id}, \texttt{puzzle\_id}
\end{description}

Tener IDs locales y globales nos permite referenciar una pieza de forma única
dentro del modelo de datos, pero al mismo tiempo mantener una representación
corta y legible para presentar al usuario. El ID global se construye
concatenando el slug del puzzle con el ID local (ver \S\ref{sec:unicidad}), y
es el que se almacena en Neo4j; el ID local es el que aparece en las
instrucciones de ensamblaje.

\subsection{Relaciones}

\begin{description}[leftmargin=4.5cm, style=sameline]
    \item[\texttt{CONTAINS}]       \texttt{(:Puzzle)} $\to$ \texttt{(:Piece)}
    \item[\texttt{HAS\_CONNECTOR}] \texttt{(:Piece)} $\to$ \texttt{(:Connector)}
    \item[\texttt{CONNECTS\_WITH}] \texttt{(:Connector)} $\leftrightarrow$ \texttt{(:Connector)} — almacenado en ambas direcciones
\end{description}

Mediante estas relaciones podemos identificar las piezas que pertenecen a cada
rompecabezas y todas las conexiones entre ellas. La relación
\texttt{CONNECTS\_WITH} se almacena en ambas direcciones deliberadamente: así
cualquier traversal de vecindad puede hacerse con un único patrón
\texttt{MATCH} sin necesidad de manejar dirección, lo que simplifica
notablemente las consultas del algoritmo.

\subsection{Diagrama del esquema}

\begin{center}
\begin{tikzpicture}[
    node distance=1.8cm and 3cm,
    box/.style={draw, rounded corners=4pt, minimum width=2.6cm,
                minimum height=0.9cm, align=center, font=\small\ttfamily},
    arr/.style={-{Stealth[length=6pt]}, thick},
]
    \node[box, fill=violet!15] (puz) {:Puzzle};
    \node[box, fill=blue!12,   right=of puz] (pc)  {:Piece};
    \node[box, fill=teal!15,   right=of pc]  (cn)  {:Connector};

    \draw[arr] (puz) -- node[above,font=\scriptsize]{CONTAINS}       (pc);
    \draw[arr] (pc)  -- node[above,font=\scriptsize]{HAS\_CONNECTOR} (cn);
    \draw[arr] (cn)  edge[loop right, looseness=6]
               node[right,font=\scriptsize]{CONNECTS\_WITH}          (cn);
\end{tikzpicture}
\end{center}

\subsection{Escalabilidad del modelo}

El modelo escala bien en ambas dimensiones relevantes. Agregar un nuevo
rompecabezas es insertar un subgrafo completamente aislado: no afecta los nodos
ni las relaciones de puzzles existentes, no requiere modificar el esquema y no
introduce ningún tipo de contención. Agregar piezas o conectores a un puzzle
existente sigue el mismo patrón: nodos nuevos enlazados al subgrafo
correspondiente. El namespacing por \texttt{puzzle\_id} garantiza que el
crecimiento horizontal entre puzzles nunca genere colisiones, y las restricciones
\texttt{UNIQUE} mantienen la integridad sin degradación de rendimiento gracias a
los índices que Neo4j crea automáticamente al definirlas.

\subsection{Unicidad global entre rompecabezas}
\label{sec:unicidad}

Neo4j comparte un único grafo de propiedades entre todos los rompecabezas
almacenados. Dado que los IDs de piezas y conectores son locales a su puzzle
(dos puzzles distintos pueden tener una pieza \texttt{P01}), es necesario
garantizar unicidad global. Para ello, los IDs de \texttt{:Piece} y
\texttt{:Connector} se construyen concatenando el slug del puzzle con el ID
local:

\[
    \texttt{id} = \texttt{"<puzzle\_id>\_\_<local\_id>"}
\]

Se crean tres restricciones \texttt{UNIQUE} (Puzzle.id, Piece.id, Connector.id)
al inicializar la base de datos, de forma que un intento de insertar dos nodos
con el mismo ID global falle a nivel de base de datos antes de llegar a la
lógica de aplicación.

% ─────────────────────────────────────────────────────────────────────────────
\section{Implementación del Modelo}

La carga de datos sigue un orden determinado: primero el nodo \texttt{:Puzzle},
luego cada \texttt{:Piece} enlazada con \texttt{CONTAINS}, luego cada
\texttt{:Connector} enlazado con \texttt{HAS\_CONNECTOR}, y finalmente las
relaciones \texttt{CONNECTS\_WITH} entre pares de conectores. Se utiliza
\texttt{MERGE} en lugar de \texttt{CREATE} para que la operación sea idempotente:
ejecutar el seed dos veces no duplica datos.

\begin{lstlisting}[language=SQL, caption={Inserción de un rompecabezas en Neo4j}]
-- Nodo Puzzle
MERGE (puz:Puzzle {id: $pid})
SET puz.name = $name, puz.description = $desc;

-- Nodo Piece
MATCH (puz:Puzzle {id: $pid})
MERGE (pc:Piece {id: $gid})
SET pc.local_id=$lid, pc.label=$label, pc.puzzle_id=$pid
MERGE (puz)-[:CONTAINS]->(pc);

-- Nodo Connector
MATCH (pc:Piece {id: $piece_gid})
MERGE (c:Connector {id: $gid})
SET c.local_id=$lid, c.label=$label, c.type=$type,
    c.piece_local_id=$plid, c.puzzle_id=$pid
MERGE (pc)-[:HAS_CONNECTOR]->(c);

-- Relacion CONNECTS_WITH (bidireccional)
MATCH (ca:Connector {id: $ga}), (cb:Connector {id: $gb})
MERGE (ca)-[:CONNECTS_WITH]->(cb)
MERGE (cb)-[:CONNECTS_WITH]->(ca);
\end{lstlisting}

El rompecabezas de prueba cargado es un \textbf{Plesiosaurio} con 10 piezas
(P01--P10), conectores etiquetados A--J y 10 conexiones que forman la columna
vertebral y extremidades del animal. Este puzzle sirve para validar tanto el
caso completo (todas las piezas disponibles) como el caso parcial (una o más
piezas marcadas como faltantes).

% ─────────────────────────────────────────────────────────────────────────────
\section{Algoritmo de Resolución}

\subsection{Descripción}

El algoritmo aplica \textbf{BFS (Breadth First Search)} sobre el grafo de
piezas, expandiendo la frontera conector por conector a partir de una pieza
inicial elegida por el usuario. Se elige BFS sobre DFS porque produce
instrucciones en orden de distancia desde la pieza inicial, lo que resulta más
natural para un ensamblador físico: primero se colocan las piezas directamente
adyacentes, luego las que están a dos pasos, y así sucesivamente.

\begin{description}[leftmargin=2.5cm]
    \item[Entrada]  Grafo $G = (P, C, E)$, pieza inicial $s \in P$, conjunto de piezas faltantes $M \subseteq P$.
    \item[Salida]   Lista ordenada de pasos de ensamblaje, cada uno con estado
                    \texttt{start}, \texttt{placed} o \texttt{missing}.
\end{description}

\subsection{Pseudocódigo}

\begin{lstlisting}[mathescape=true, caption={BFS de ensamblaje}]
visited = {s};  queue = [s]
emitir paso START: "Coloca la pieza s como punto de partida"

mientras queue no este vacio:
    u = queue.pop_left()
    para cada conector $c_u$ de u:
        $c_v$ = vecino($c_u$)          // CONNECTS_WITH en Neo4j
        v  = pieza($c_v$)
        si v en M (faltante) y no reportada:
            emitir paso MISSING
            continuar
        si v en visited: continuar
        visited.add(v);  queue.append(v)
        emitir paso PLACED:
            "Toma la pieza v, conecta su $c_v$.tipo '$c_v$.label'
             al $c_u$.tipo '$c_u$.label' de la pieza u"

para cada p en P - visited - M:
    emitir paso MISSING: "pieza p disponible pero inalcanzable"
\end{lstlisting}

La complejidad del algoritmo es $\mathcal{O}(n + m)$ donde $n = |P|$ es el
número de piezas y $m = |E|$ el número de conexiones. Cada pieza se visita a lo
sumo una vez y cada conexión se examina a lo sumo dos veces (una desde cada
extremo), por lo que el algoritmo escala linealmente con el tamaño del
rompecabezas.

\subsection{Uso de Neo4j en el algoritmo}

Al cargar el rompecabezas, se ejecuta una única consulta Cypher que materializa
el mapa completo de vecinos (\texttt{connector\_id} $\to$ \texttt{connector\_id})
en memoria:

\begin{lstlisting}[language=SQL]
MATCH (:Puzzle {id: $pid})-[:CONTAINS]->(:Piece)
      -[:HAS_CONNECTOR]->(ca:Connector)
      -[:CONNECTS_WITH]->(cb:Connector)
WHERE ca.puzzle_id = $pid AND cb.puzzle_id = $pid
  AND ca.local_id < cb.local_id
RETURN ca.local_id AS a, cb.local_id AS b;
\end{lstlisting}

El BFS opera sobre este mapa en memoria; la base de datos provee la
representación persistente y compartida. Dado que la cantidad de datos por
rompecabezas es pequeña (decenas o cientos de piezas), cargar el grafo completo
de una sola vez es más eficiente que consultar Neo4j por cada paso del BFS,
evitando la latencia de red en cada iteración.

% ─────────────────────────────────────────────────────────────────────────────
\section{Manejo de Piezas Faltantes}

Una de las funcionalidades clave del sistema es continuar generando
instrucciones útiles incluso cuando algunas piezas no están disponibles. Cuando
durante el BFS se encuentra una pieza $v \in M$:

\begin{enumerate}
    \item Se registra un paso \texttt{MISSING} indicando qué conector debería
          unirse y de qué pieza, para que el ensamblador sepa qué falta.
    \item No se encola $v$, bloqueando esa rama de expansión y evitando
          instrucciones que no se pueden ejecutar.
    \item Las piezas alcanzables por rutas alternativas (que no pasen por $v$)
          se ensamblan con normalidad, aprovechando al máximo las piezas disponibles.
    \item Al finalizar, las piezas disponibles pero inalcanzables debido al
          bloqueo se reportan explícitamente, dando al ensamblador una imagen
          completa del estado del rompecabezas.
\end{enumerate}

Esto refleja el comportamiento físico deseado: el ensamblador arma todo lo
posible, recibe una notificación clara de qué sección no puede completarse y
por qué, y puede retomar el armado si la pieza faltante aparece posteriormente.

\end{document}
